% Part: turing-machines 
% Chapter: undecidability 
% Section: unsolvability-decision-problem

\documentclass[../../../include/open-logic-section]{subfiles}

\begin{document}

\olfileid[pt]{tur}{und}{uns} 
\olsection{O Problema da Decisão é Insolúvel}

\begin{thm}
\ollabel{thm:decision-prob}
O problema da decisão é insolúvel: Não há máquina de Turing~$D$ tal
que, quando iniciada em uma fita que contém como entrada
!!a{sentence}~$!B$ da lógica de primeira ordem, $D$~eventualmente para
e produz~$1$ se, e somente se, $!B$ é válida, e $0$ caso contrário.
\end{thm}

\begin{proof}
Suponha que o problema da decisão fosse solúvel, isto é, suponha que
houvesse uma máquina de Turing~$D$. Então poderíamos resolver o
problema da parada como segue. Construímos uma máquina de Turing~$E$
que, dado como entrada o número~$e$ da máquina de Turing~$M_e$ e a
entrada~$w$, computa a !!{sentence} correspondente~$!T(M_e, w) \lif
!E(M_e, w)$ e para, examinando o quadrado mais à esquerda na fita.  A
máquina $E \concat D$ então, dada a entrada $e$ e $w$, primeiro
computaria~$!T(M_e, w) \lif !E(M_e, w)$ e depois executaria a máquina
do problema da decisão~$D$ nessa entrada.  $D$ para com saída~$1$ se,
e somente se, $!T(M_e, w) \lif !E(M_e, w)$ é válida, e produz~$0$ caso
contrário. Por \olref[ver]{lem:halt-if-valid} e
\olref[ver]{lem:valid-if-halt}, $!T(M_e, w) \lif !E(M_e, w)$ é válida
se, e somente se, $M_e$ para na entrada~$w$. Assim, $E\concat D$, dada
a entrada $e$ e $w$, para com saída~$1$ se, e somente se, $M_e$ para na
entrada~$w$, e para com saída~$0$ caso contrário. Em outras palavras,
$E \concat D$ resolveria o problema da parada.  Mas sabemos, por
\olref[hal]{thm:halting-problem}, que não pode existir tal máquina de
Turing.
\end{proof}

\begin{cor}\ollabel{cor:undecidable-sat}%
É indecidível se uma !!{sentence} arbitrária da lógica de primeira ordem é satisfatível.
\end{cor}

\begin{proof}
  Suponha que a satisfatibilidade fosse decidível por uma máquina de
  Turing~$S$. Então poderíamos resolver o problema da decisão como
  segue: Dada como entrada !!a{sentence}~$B$, mova $!B$ um quadrado
  para a direita. Volte ao quadrado~$1$ e escreva o símbolo~$\lnot$.

  Agora execute a máquina de Turing~$S$. Ela eventualmente para com
  saída $1$ (se $\lnot !B$ é satisfatível) ou~$0$ (se $\lnot !B$ é
  insatisfatível) na fita. Se há um~$\TMstroke$ no quadrado~$1$,
  apague-o; se o quadrado~$1$ está vazio, escreva um~$\TMstroke$ e
  então pare.

  Essa máquina de Turing sempre para, e sua saída é~$1$ se, e somente
  se, $\lnot !B$ é insatisfatível, e $0$~caso contrário. Como $!B$ é
  válida se, e somente se, $\lnot !B$~é insatisfatível, a máquina
  produz~$1$ se, e somente se, $!B$ é válida, e $0$~caso contrário,
  isto é, ela resolveria o problema da decisão.
\end{proof}

\begin{explain}
Portanto, não há máquina de Turing que sempre dê uma resposta ``sim''
ou ``não'' correta à pergunta ``$!B$ é uma !!{sentence} válida da
lógica de primeira ordem?'' No entanto, \emph{há} uma máquina de
Turing que sempre dá uma resposta ``sim'' correta---mas que
simplesmente não para se a resposta é ``não''. Isso decorre do teorema
da correção e completude da lógica de primeira ordem, e do fato de que
!!{derivation}s podem ser efetivamente enumeradas.
\end{explain}

\begin{thm}
  \ollabel{thm:valid-ce}%
  A validade de !!{sentence}s de primeira ordem é semidecidível: Há
  uma máquina de Turing~$E$ que, quando iniciada em uma fita que
  contém como entrada !!a{sentence}~$!B$ da lógica de primeira ordem,
  $E$~eventualmente para e produz~$1$ se, e somente se, $!B$ é válida,
  mas não para caso contrário.
\end{thm}

\begin{proof}
  Todas as !!{derivation}s possíveis da lógica de primeira ordem podem
  ser geradas, uma após a outra, por um algoritmo efetivo.  A
  máquina~$E$ faz isso e, quando encontra !!a{derivation} que mostra
  que $\Proves !B$, ela para com saída~$1$. Pelo teorema da correção,
  se $E$ para com saída~$1$, é porque~$\Entails !B$. Pelo teorema da
  completude, se $\Entails !B$, há !!a{derivation} que mostra
  que~$\Proves !B$. Como $E$ gera sistematicamente todas as
  !!{derivation}s possíveis, ela eventualmente encontrará uma que
  mostra~$\Proves !B$, de modo que eventualmente parará com saída~$1$.
\end{proof}

\end{document}
